\subsection{Visualizer}
\label{sec:visualizer-methods}
% Order follows how a user works with the app:
% load -> see the whole recording -> zoom in -> find activity -> check the pieces.
% Authors: Keito = overview, page layout, overview window, activity heatmap;
%          Kazuki = coarse-first loading, detail window (incl. LTTB), event list.

% ---- Overview (Keito) ----
% TODO (Keito): purpose and requirements, workflow, why Streamlit + Plotly,
% architecture (UI in app.py/panels/state.py, Streamlit-free viewport/plots/signal_io/detectors;
% same core/ code as the command line).
Large Marginum CSV files are several gigabytes at \SI{50}{\kilo\hertz}, so they cannot be
loaded and plotted as a single dense series in a browser. We built an interactive
Streamlit application that loads a recording, shows a coarse overview, and lets the user
inspect selected time windows in detail.

\subsubsection{Page layout}
\label{sec:vis-layout}
% TODO (Keito): three columns (load and settings / plots / results), one module per area
% in panels/, state kept in state.py, pluggable detectors (common Detection table),
% each column scrolls on its own. Figure: images/viewer.png.

\subsubsection{Coarse-first loading}
\label{sec:vis-loading}
Opening a file first builds a lightweight preview (and a row index) so the overview can
appear quickly. In the background the recording is reduced to a \SI{1}{\kilo\hertz} summary
(mean and envelope). After that summary is ready, detection and accurate detail views use
it; until then the UI remains in a preview state. Index and summary are cached so later
opens of the same file are much faster.

\subsubsection{Overview window}
\label{sec:vis-overview}
% TODO (Keito): full recording with a marker per piece; the frame selects what the detail
% window shows; drag a box / move the slider handles / window size; limited to 3 minutes
% (reason: refer to Section~\ref{sec:vis-detail}).

\subsubsection{Detail window}
\label{sec:vis-detail}
% TODO (Kazuki): exact samples of both channels in the frame, pieces marked,
% zoom / pan / default / fit all buttons.

\subsubsection{Detail downsampling with LTTB}
\label{sec:vis-lttb}
Even a short detail window can contain tens or hundreds of thousands of samples. Plotting
every sample is slow and visually cluttered. We therefore downsample the detail traces for
display with the Largest-Triangle-Three-Buckets (LTTB) algorithm of Steinarsson
\cite{steinarsson2013} (Algorithm~4.2), implemented in the visualizer as \texttt{lttb.py}.

LTTB partitions the series into buckets and, in each bucket, keeps the point that forms the
largest triangle with the previously selected point and the average of the next bucket.
Compared with taking every $n$th sample (stride), LTTB better preserves perceptually
important peaks and troughs at the same point budget, which matters when reviewing
fluorescence spikes from aspirated tissue.

\subsubsection{Activity heatmap}
\label{sec:vis-heatmap}
% TODO (Keito): recording folded into rows (row length 10 s -- 5 min, cells per row),
% colour = signal activity or number of pieces; click a cell to move the window there
% (window jump described once, here or in Section~\ref{sec:vis-events}).

\subsubsection{Event list}
\label{sec:vis-events}
% TODO (Kazuki): accepted pieces largest first, amount in view, "Show in plot",
% selected event and detection summary.
